\Needspace{8\baselineskip}\originalchapter{带电粒子轨道与绝热不变量}\label{ch:orbits}
\originalsection{均匀场与回旋}
Lorentz 方程 $m\dot{\mathbf v}=q(\mathbf E+\mathbf v\times\mathbf B)$ 中, 磁力与速度垂直, 单独不做功. 取 $\mathbf B=B\hat z$, 无电场时横向运动为频率 $\Omega=qB/m$ 的回旋, Larmor 半径 $r_L=mv_\perp/(|q|B)$, 平行运动匀速. $\Omega$ 保留电荷符号, 描述旋转方向; 正的频率幅值写 $|\Omega|$.

均匀垂直电场下, 找一个速度 $\mathbf v_E$ 使 $\mathbf E+\mathbf v_E\times\mathbf B=0$, 得 $\mathbf v_E=\mathbf E\times\mathbf B/B^2$. 剩余速度仍回旋. 此漂移与质量、电荷无关, 电子和离子的共同移动本身不产生净流电流. 平行电场则持续加速, 不能用这一个漂移消去.

\originalsection{导引中心漂移}
场在一回旋尺度上慢变, $r_L/L_B\ll1$, 时间变化也慢 $|\dot B|/(B|\Omega|)\ll1$. 平均一周期后, 额外缓慢运动可用导引中心描述. 定义磁矩 $\mu=mv_\perp^2/(2B)$, 梯度漂移与曲率漂移为
\begin{equation}\label{eq:drifts}
 \mathbf v_{\nabla B}=\frac{\mu}{qB}\mathbf b\times\nabla B,
 \qquad\mathbf v_c=\frac{mv_\parallel^2}{qB}\mathbf b\times\boldsymbol\kappa,
 \quad\boldsymbol\kappa=(\mathbf b\cdot\nabla)\mathbf b.
\end{equation}
梯度漂移来自回旋两侧半径和速度的微差; 曲率漂移来自沿弯曲磁线运动的惯性力. 二者依 $1/q$, 正负粒子往往反向, 可形成电流. 一般非磁力 $\mathbf F$ 给 $\mathbf v_F=\mathbf F\times\mathbf B/(qB^2)$.

电场缓变还产生极化漂移 $\mathbf v_p=m\dot{\mathbf E}_\perp/(qB^2)$. 离子与电子质量差使净极化电流以离子贡献为主. 这些式在 B趋零或 $E/B$ 接近 c 的区域失去非相对论与绝热控制.

\originalsection{磁镜与损失锥}
无电场且绝热时 $W=mv^2/2$ 与 $\mu$ 同时守恒. 沿磁线到高B处, $W_\perp=\mu B$ 增加而 $W_\parallel=W-\mu B$ 减少. 当 $B=B_{
\rm turn}=W/\mu$ 时平行速度零并反射. 中心场 $B_0$、喉场 $B_1$, 初始俯仰角 $\alpha$ 满足 $\sin^2\alpha\ge B_0/B_1$ 才能被困. 其他角位于两端损失锥.

\noindent\begin{minipage}{\linewidth}
\begin{center}\begin{tikzpicture}[x=1cm,y=1cm]
\draw[->](-2,0)--(2,0)node[right]{$v_\parallel$};\draw[->](0,0)--(0,1.8)node[above]{$v_\perp$};
\draw[dashed](0,0)--(1.7,1.0);\draw[dashed](0,0)--(-1.7,1.0);\node at(0,1.2){被困};\node at(1.6,.3){损失};\node at(-1.6,.3){损失};
\end{tikzpicture}\end{center}
图为速度空间示意, 损失锥角由镜比决定, 不同分布的被困比例还需角积分. 碰撞可将被困粒子散入锥内, 因而磁镜不能由单粒子不变量保证无限时间约束.
\end{minipage}

\originalsection{能量、磁矩与时变场}
导引中心哈密顿量近似 $W=m v_\parallel^2/2+\mu B+q\phi$. 静场沿磁线有镜力 $m\dot v_\parallel=-\mu\nabla_\parallel B+qE_\parallel$. 缓慢增加 B 会对回旋粒子作 betatron 加热, 此时 \(W_\perp\propto B\), 来源是 Faraday 感应电场, 不是磁力直接做功.

对时间变化的电磁场, 只计算静电势能并不能保证机械加电势能守恒. 能量可从外部电源流入. 不变量有明确慢变条件; 回旋轨道穿越场零点、非绝热电场或强碰撞时可能破坏它们.

\originalsection{电场主导的相对论运动}
相对论动量 $\mathbf p=\gamma m\mathbf v$ 满足 $\dot{\mathbf p}=q(\mathbf E+\mathbf v\times\mathbf B)$. 当互相垂直的 E、B 有 $E<cB$, 可以 boost 到 E'=0 系, 看到回旋加整体漂移. 当 E>cB, 则可用速度 $\mathbf u=c^2\mathbf E\times\mathbf B/E^2$ boost 到 B'=0 系, $|u|<c$, 粒子受纯电场加速, 轨道不闭合. $E/B$ 此时不能作为实际粒子漂移速度, 否则会错误超过 c. 两场不变量 $B^2-E^2/c^2$ 与 $\mathbf E\cdot\mathbf B$ 决定能否消去其中一场.
